\chapter{Implémentation}
\label{chap:implementation}

\section{Environnement et chaîne de compilation}

Le projet est écrit en C~11 et compilé avec \code{gcc} (ou
\code{clang}) en utilisant un \file{Makefile} GNU classique. Les
options de compilation par défaut sont~:

\begin{minted}{make}
CFLAGS ?= -std=c11 -O2 -Wall -Wextra -Wpedantic
\end{minted}

Les drapeaux GTK4 sont obtenus dynamiquement via \code{pkg-config}~:

\begin{minted}{make}
GTK_CFLAGS := $(shell $(PKG_CONFIG) --cflags gtk4)
GTK_LIBS   := $(shell $(PKG_CONFIG) --libs gtk4)
\end{minted}

L'arborescence du \emph{build} est~:

\begin{itemize}
  \item \file{build/obj/} reçoit les fichiers objets~;
  \item \file{build/bin/aquarium} est le binaire final~;
  \item \file{vendor/tomlc99/} contient \file{toml.h}, \file{toml.c}
    et la licence MIT du parseur.
\end{itemize}

Une cible \code{install-conf} du \file{Makefile} copie
\file{aquarium.conf} dans \file{\textasciitilde/.config/gtkaqua/}
pour faciliter l'usage du \emph{fallback} système.

\section{Pipeline d'assets}

Les sprites bruts vivent dans \file{Fish/} (formats \file{webp}
ou \file{png}). La cible \code{make assets} les normalise vers
\file{assets/} en garantissant qu'ils sont tous orientés vers la
gauche. Pour les fichiers d'origine déjà gauche, la commande se
limite à un redimensionnement~; pour les fichiers déjà droite, un
flip horizontal est appliqué via \code{ffmpeg}~\cite{ffmpeg}~:

\begin{minted}{make}
ffmpeg -i $(FISH_DIR)/shark.webp -vf "scale=100:-1" \
        -frames:v 1 -y $(ASSET_DIR)/shark.png
ffmpeg -i $(FISH_DIR)/fish1Texture.png -vf "hflip" \
        -frames:v 1 -y $(ASSET_DIR)/goldfish.png
\end{minted}

Cette discipline d'orientation est documentée dans le
\file{README.md}~: au moment du runtime, le code suppose que tous
les sprites pointent à gauche. La rotation de chaque sprite est
ensuite produite par une transformation CSS GTK4 sur le widget,
sans réchantillonnage de l'image.

\section{Algorithmes clés}

Cette section présente les fonctions centrales du moteur de
simulation. Pour des raisons de place, on extrait les passages
significatifs~; le code complet est disponible en annexe et dans
le dépôt.

\subsection{Boucle de tick}

La fonction \func{world\_tick} orchestre une itération complète
de la simulation. Elle est listée dans le code~\ref{lst:world-tick}.

\begin{listing}[h!]
\inputminted[firstline=298, lastline=313]{c}{../src/world.c}
\caption{La fonction \func{world\_tick} (extrait de \file{src/world.c}).}
\label{lst:world-tick}
\end{listing}

L'ordre est important~: déplacement physique, puis interactions
(repas, famine), puis mise à jour visuelle, puis gestion des
respawns. Cette séquence évite les races logiques entre événements
du même tick.

\subsection{Forces de Reynolds}

Le calcul des trois forces de banc (séparation, alignement,
cohésion) constitue le cœur algorithmique du moteur. Le code
parcourt tous les voisins de la \emph{même} espèce dans le rayon
correspondant à chaque force, accumule les contributions, puis
construit un vecteur de \emph{steering} pondéré par les coefficients
de la \struct{SpeciesConfig}.

\begin{listing}[ht]
\inputminted[firstline=56, lastline=125]{c}{../src/entity.c}
\caption{Calcul des trois forces de banc (extrait de \file{src/entity.c}).}
\label{lst:force-boids}
\end{listing}

Trois remarques sur ce code~:

\begin{itemize}
  \item Les distances sont comparées au carré
    (\code{dist\_sq < radius * radius}), ce qui évite un
    \code{sqrt} par voisin et constitue une optimisation
    classique.
  \item L'\emph{enveloppement torique} du monde est pris en
    compte via \func{vec2\_delta\_wrap}, qui choisit le delta
    minimum entre deux points à travers les bords.
  \item Le vecteur de séparation est pondéré par
    $1/\|\vec{r}\|^2$, ce qui rend la force d'autant plus forte
    que les voisins sont proches~: c'est la formulation classique
    de Reynolds.
\end{itemize}

\subsection{Force de fuite avec adrénaline}

La fuite face à un prédateur est implémentée par
\func{force\_flee}. Lorsqu'au moins un prédateur entre dans le
\code{fear\_radius} de l'agent, sa vitesse maximale et sa force
maximale sont multipliées par le facteur \code{g\_adrenaline\_boost}
(2.5 par défaut), ce qui rend la fuite beaucoup plus rapide que
le déplacement nominal.

\begin{listing}[ht]
\inputminted[firstline=127, lastline=173]{c}{../src/entity.c}
\caption{Force de fuite avec boost d'adrénaline (extrait de \file{src/entity.c}).}
\label{lst:force-flee}
\end{listing}

Le test \func{EATS} sur la ligne pivot de cette fonction reflète
exactement la définition donnée dans \file{species.h}~:

\begin{minted}{c}
#define EATS(predator_cfg, prey_index) \
        ((predator_cfg).prey_mask & (1 << (prey_index)))
\end{minted}

\subsection{Alimentation par la bouche}

Plutôt que de tester une collision centre-à-centre entre
prédateur et proie, le projet utilise une \emph{mouth-based
collision}~: le point d'attaque est dérivé du cap du prédateur,
selon le ratio \code{mouth\_forward\_ratio}.

\begin{listing}[h!]
\inputminted[firstline=33, lastline=41]{c}{../src/entity.c}
\caption{Calcul du point de bouche d'un prédateur.}
\label{lst:mouth}
\end{listing}

Ce point est ensuite comparé à la position centrale de chaque
proie compatible. La proie est mangée si la distance enveloppée
est inférieure à \code{g\_eat\_radius + prey\_hit}, où
\code{prey\_hit} dépend de la taille du sprite de la proie. La
formule pour \code{prey\_hit} accorde un rayon de hit
proportionnel à la plus petite dimension du sprite, mais avec un
plancher de 8~pixels pour ne pas pénaliser les sprites
microscopiques.

\begin{listing}[h!]
\inputminted[firstline=196, lastline=237]{c}{../src/world.c}
\caption{Détection des repas (extrait de \file{src/world.c}).}
\label{lst:eating}
\end{listing}

\subsection{Cycle de vie : despawn et respawn}

\func{world\_despawn} fait passer une entité de \code{ALIVE} à
\code{DYING}, décrémente le compteur d'espèce, arme le timer de
flash de mort, et programme la prochaine fenêtre de respawn.

\begin{listing}[h!]
\inputminted[firstline=116, lastline=138]{c}{../src/world.c}
\caption{Despawn d'une entité (extrait de \file{src/world.c}).}
\label{lst:despawn}
\end{listing}

Symétriquement, \func{world\_handle\_respawns} parcourt les
espèces dont la population est inférieure à \code{flock\_size},
attend la fenêtre \code{respawn\_ready}, puis recycle une entité
\code{INACTIVE} en lui assignant une nouvelle position et une
nouvelle vitesse aléatoires.

\begin{listing}[h!]
\inputminted[firstline=252, lastline=296]{c}{../src/world.c}
\caption{Logique de respawn (extrait de \file{src/world.c}).}
\label{lst:respawn}
\end{listing}

\subsection{Famine}

La famine est délibérément simple~: un prédateur qui n'a pas
mangé depuis \code{g\_starvation\_seconds} secondes est despawnée.
Pour les espèces non-prédatrices, le champ \code{last\_eaten} est
initialisé à l'instant courant et n'est jamais mis à jour~; à
moins que \code{g\_starvation\_seconds} ne soit anormalement
court, ces espèces ne meurent donc pas de faim. Cela laisse
intacte la mécanique de respawn pour les proies, qui meurent
exclusivement par prédation.

\begin{listing}[h!]
\inputminted[firstline=239, lastline=250]{c}{../src/world.c}
\caption{Vérification de la famine.}
\label{lst:lifespan}
\end{listing}

\section{Système de configuration TOML}

\subsection{Structure du fichier}

Le fichier \file{aquarium.conf} comporte une section
\code{[global]} pour les constantes du moteur, et autant de
blocs \code{[[species]]} qu'il y a d'espèces. La syntaxe TOML est
strictement définie par la spécification 1.0.0~\cite{tomlspec}.

\begin{listing}[h!]
\begin{minted}[fontsize=\footnotesize, frame=lines]{toml}
[global]
eat_radius = 10.0
adrenaline_boost = 2.5
rotation_smooth = 0.1
death_flash_duration = 0.22
starvation_seconds = 120.0
start_width = 1280
start_height = 720
start_fullscreen = false
tick_ms = 16
assets_dir = "assets/"

[[species]]
id = "goldfish"
name = "Goldfish"
sprite = "assets/goldfish.png"
sprite_width = 46
sprite_height = 23
mouth_forward_ratio = 0.40
max_speed = 85.0
min_speed = 30.0
max_force = 42.0
separation_radius = 24.0
alignment_radius = 120.0
cohesion_radius = 100.0
separation_weight = 1.6
alignment_weight = 1.35
cohesion_weight = 1.45
fear_radius = 190.0
fear_weight = 2.5
hunt_radius = 0.0
hunt_weight = 0.0
avoid_same_radius = 0.0
avoid_same_weight = 0.0
prey = []
lifespan = 120.0
flock_size = 45
max_population = 60
respawn_delay = 1.5
is_apex = false
\end{minted}
\caption{Extrait de \file{aquarium.conf} : section globale et bloc \code{goldfish}.}
\label{lst:conf-extract}
\end{listing}

\subsection{Parseur \code{tomlc99}}

La bibliothèque \code{tomlc99}~\cite{tomlc99} expose une API en
deux étapes~: \func{toml\_parse\_file} retourne un
\code{toml\_table\_t *} si le fichier est syntaxiquement valide,
puis des accesseurs typés (\func{toml\_int\_in},
\func{toml\_double\_in}, \func{toml\_bool\_in},
\func{toml\_string\_in}) lisent les valeurs par clé. Pour les
tableaux de tables, \func{toml\_array\_in} récupère le tableau,
puis \func{toml\_table\_at} accède aux sections individuelles.

Un point d'attention~: les valeurs \code{string} retournées par
\code{tomlc99} sont allouées dynamiquement et doivent être
libérées avec \code{free()} (et non pas \func{toml\_free},
qui sert à libérer la table racine). Cette subtilité a été
identifiée à la première compilation et corrigée.

\subsection{Résolution des références \code{prey}}

Le champ \code{prey} de chaque espèce est un tableau d'identifiants.
Comme les indices de la table \code{species} ne sont connus
qu'après lecture de toutes les espèces, la résolution se fait en
\textbf{deux passes}~:

\begin{enumerate}
  \item \textbf{Passe~1.} Lecture des champs scalaires de chaque
    \code{[[species]]}~: identifiant, nom, sprite, dimensions,
    paramètres comportementaux. Le \code{prey\_mask} est laissé à
    zéro.
  \item \textbf{Passe~2.} Pour chaque espèce, parcours du tableau
    \code{prey}, recherche linéaire de l'identifiant dans la table
    déjà construite, et activation du bit correspondant dans
    \code{prey\_mask}. Une référence inconnue déclenche une erreur
    de chargement nommant l'espèce coupable et la proie absente.
\end{enumerate}

Cette double passe est exécutée pendant que la \code{toml\_table\_t}
racine est encore vivante en mémoire, ce qui permet de simplement
re-parcourir l'arbre TOML sans avoir à mémoriser les chaînes brutes
entre les deux passes.

\subsection{Chemins de recherche et fallback}

\func{conf\_load} essaie successivement plusieurs chemins, dans
l'ordre suivant~:

\begin{enumerate}
  \item \code{override\_path} explicite passé en argument
    (réservé à un futur \code{--config}).
  \item \file{./aquarium.conf} (répertoire courant).
  \item \file{\$XDG\_CONFIG\_HOME/gtkaqua/aquarium.conf}, ou à
    défaut \file{\textasciitilde/.config/gtkaqua/aquarium.conf}.
  \item \file{/etc/gtkaqua/aquarium.conf} (configuration système).
\end{enumerate}

Le premier chemin existant et lisible (\func{stat} qui renvoie
un fichier régulier) est utilisé. Si aucun n'est trouvé, un
message d'erreur précis liste les chemins essayés.

\subsection{Variables globales et runtime}

Les paramètres globaux issus de la configuration sont publiés
sous forme de variables \code{extern} déclarées dans
\file{include/config.h} et définies dans \file{src/conf.c}~:

\begin{minted}{c}
extern double g_eat_radius;
extern double g_adrenaline_boost;
extern double g_rotation_smooth;
extern double g_death_flash_duration;
extern double g_starvation_seconds;
extern int    g_start_width;
extern int    g_start_height;
extern int    g_start_fullscreen;
extern int    g_tick_ms;
extern const char *g_assets_dir;
\end{minted}

Ce choix simplifie la consommation~: \file{src/world.c} et
\file{src/entity.c} référencent ces variables sans avoir besoin
de connaître la \struct{AquariumConfig}. Le couplage est
unidirectionnel~: \code{conf} écrit, tout le reste lit.

\subsection{Rechargement à chaud}

Le bouton \emph{Reload config} de la barre latérale appelle un
gestionnaire qui~:

\begin{enumerate}
  \item Construit un \struct{AquariumConfig} temporaire et appelle
    \func{conf\_load} dessus.
  \item Si l'opération échoue, affiche une boîte de dialogue
    d'erreur GTK avec le message d'erreur retourné par
    \func{conf\_error}, et abandonne sans toucher à l'état
    courant.
  \item Si l'opération réussit, copie le tableau d'espèces dans le
    registre runtime via \func{species\_init}, puis vide le monde
    (\func{world\_clear}), reconstruit dynamiquement l'interface
    des espèces (cases à cocher du \emph{pool}, combobox
    \emph{Add precise}), et appelle \func{world\_populate} pour
    repeupler avec la nouvelle configuration.
\end{enumerate}

Cette opération préserve la fenêtre, la pause, et tous les
contrôles non liés aux espèces.

\section{Rendu GTK4}

\subsection{Hiérarchie des widgets}

L'application repose sur une hiérarchie GTK simple~:

\begin{itemize}
  \item \code{GtkApplicationWindow}.
  \item \code{GtkBox} horizontal racine.
  \item À gauche, un \code{GtkScrolledWindow} contenant la barre
    latérale (\code{GtkBox} vertical) avec les contrôles.
  \item À droite, un \code{GtkOverlay} qui empile la vidéo de
    fond (\code{GtkPicture} avec \code{GtkMediaFile}) et un
    \code{GtkFixed} servant de toile pour les sprites.
\end{itemize}

\code{GtkFixed} est un conteneur qui place ses enfants en
coordonnées absolues. À chaque tick, \func{entity\_apply\_visuals}
appelle \func{gtk\_fixed\_move} pour mettre à jour la position
de chaque sprite.

\subsection{Rotation par CSS}

La rotation des sprites est confiée à GTK via une transformation
CSS appliquée à chaque widget. Chaque entité possède un nom de
classe CSS unique (\code{entity-N}) et un \code{GtkCssProvider}
dédié. À chaque tick, le CSS est régénéré~:

\begin{minted}{c}
snprintf(css, sizeof(css),
         ".%s { transform: rotate(%.1fdeg); opacity: 1.0; }",
         e->css_class, e->angle);
gtk_css_provider_load_from_string(e->css_provider, css);
\end{minted}

L'angle est lissé entre l'angle courant et la cible
(\code{heading\_deg\_from\_vel}) selon le facteur
\code{g\_rotation\_smooth}~:

\begin{minted}{c}
double diff = target_deg - e->angle;
while (diff > 180.0)  diff -= 360.0;
while (diff < -180.0) diff += 360.0;
e->angle += diff * g_rotation_smooth;
\end{minted}

Cela évite les changements brusques de direction visuelle.

\subsection{Effet de mort}

L'état \code{DYING} active un effet visuel pulsé combinant
rotation, mise à l'échelle et variation d'opacité~:

\begin{minted}{c}
double pulse   = fmod(e->death_timer * 24.0, 2.0) < 1.0 ? 1.0 : 0.0;
double scale   = 1.15 + (0.35 * pulse);
double opacity = 0.35 + (0.55 * pulse);
snprintf(css, sizeof(css),
         ".%s { transform: rotate(%.1fdeg) scale(%.2f); "
         "opacity: %.2f; }",
         e->css_class, e->angle, scale, opacity);
e->death_timer -= w->dt;
if (e->death_timer <= 0.0) {
    e->state = ENTITY_INACTIVE;
    gtk_widget_set_visible(e->widget, FALSE);
}
\end{minted}

Le calcul \code{fmod(e->death\_timer * 24.0, 2.0) < 1.0} produit
un battement à 12~Hz qui alterne entre deux états visuels,
donnant un effet de \emph{flash} caractéristique pendant
\code{g\_death\_flash\_duration} secondes.

\section{Récapitulatif modulaire}

Au final, le projet livre une base de code de l'ordre de 1500
lignes de C (hors \emph{vendor}), répartie selon le
tableau~\ref{tab:loc}.

\begin{table}[h!]
\centering
\footnotesize
\begin{tabular}{@{}lr@{}}
\toprule
\textbf{Fichier} & \textbf{Lignes (approx.)} \\
\midrule
\file{src/main.c}     & 560 \\
\file{src/world.c}    & 335 \\
\file{src/entity.c}   & 380 \\
\file{src/species.c}  &  35 \\
\file{src/conf.c}     & 305 \\
\file{include/*.h}    & 200 \\
\file{aquarium.conf}  & 240 \\
\midrule
\textbf{Total projet} & $\approx$ 2050 \\
\bottomrule
\end{tabular}
\caption{Volume de code par module.}
\label{tab:loc}
\end{table}

Cette répartition reflète la philosophie du projet~: un coût
fixe modeste pour chaque module sauf \file{main.c} (qui porte
toute l'interface utilisateur GTK) et \file{entity.c} (qui porte
les algorithmes physiques et de pilotage).
